\documentclass[14pt,a4paper]{extarticle}
\usepackage[utf8]{inputenc}
\usepackage[T1]{fontenc}
\usepackage[russian]{babel}
\usepackage{geometry}
\geometry{top=2cm, bottom=2cm, left=2.5cm, right=2cm}
\usepackage{graphicx}
\usepackage{hyperref}
\usepackage{listings}
\usepackage{xcolor}
\usepackage{enumitem}
\usepackage{titlesec}
\usepackage{setspace}
\usepackage{tikz}
\usepackage{pgf-umlcd}
\onehalfspacing

\lstset{
    basicstyle=\footnotesize\ttfamily,
    breaklines=true,
    frame=single,
    numbers=left,
    numberstyle=\tiny,
    tabsize=2,
    captionpos=b,
    showstringspaces=false,
    language=Java,
    keywordstyle=\color{blue},
    stringstyle=\color{red},
    commentstyle=\color{gray},
}

\hypersetup{
    colorlinks=true,
    linkcolor=black,
    urlcolor=blue,
}

% UML styling
\renewcommand{\umltextcolor}{black}
\renewcommand{\umldrawcolor}{black}
\renewcommand{\umlfillcolor}{white}

\begin{document}

\begin{titlepage}
    \centering
    {\scshape\large Министерство науки и высшего образования Российской Федерации\par}
    \vspace{0.3cm}
    {\scshape\large Университет ИТМО\par}
    \vspace{0.3cm}
    {\scshape\large Факультет программной инженерии и компьютерной техники\par}
    \vspace{2cm}

    {\Large\bfseries Информационные системы\par}
    \vspace{0.5cm}
    {\Large\bfseries Лабораторная работа №~1\par}
    \vspace{0.5cm}
    {\large Разработка информационной системы для управления объектами\par}
    \vspace{0.3cm}
    {\large класса \texttt{HumanBeing}\par}

    \vspace{3cm}
    \begin{flushright}
        \begin{tabular}{rl}
            \textbf{Исполнитель:} & Ясаков~А. \\[0.3cm]
            \textbf{Преподаватель:} & Кобик~Н.\,А. \\
        \end{tabular}
    \end{flushright}

    \vfill
    {\large Санкт-Петербург\par}
    {\large 2026\par}
\end{titlepage}

\tableofcontents
\newpage

%============================================================
\section{Текст задания}

Реализовать информационную систему, которая позволяет взаимодействовать с объектами класса \texttt{HumanBeing}.

\subsection{Описание классов предметной области}

\begin{lstlisting}[language=Java]
public class HumanBeing {
    private Integer id;          // not null, >0, unique, auto-generated
    private String name;         // not null, not empty
    private Coordinates coordinates; // not null
    private java.time.LocalDateTime creationDate; // not null, auto-generated
    private boolean realHero;
    private boolean hasToothpick;
    private Car car;             // nullable
    private Mood mood;           // nullable
    private float impactSpeed;   // max 738
    private WeaponType weaponType; // not null
}

public class Coordinates {
    private Float x;             // not null, max 256
    private float y;             // > -533
}

public class Car {
    private String name;         // nullable
}

public enum Mood {
    SADNESS, APATHY, CALM, FRENZY;
}

public enum WeaponType {
    PISTOL, KNIFE, BAT;
}
\end{lstlisting}

\subsection{Требования к функционалу}

\begin{enumerate}[leftmargin=*]
    \item CRUD-операции в отдельных интерфейсах: создание, получение по ID (с выводом связанных объектов), обновление, удаление.
    \item При создании — возможность связать объект с существующими \texttt{Coordinates} и \texttt{Car}.
    \item Все операции на серверной стороне, синхронизация с БД.
    \item Главный экран: таблица со всеми атрибутами, пагинация.
    \item Фильтрация и сортировка по строковым колонкам (только полное совпадение).
    \item Переход к редактированию из таблицы и из области детального просмотра.
    \item Real-time синхронизация между клиентами при добавлении/изменении/удалении.
    \item При удалении — проверка связанных объектов, отмена с сообщением при наличии связей.
    \item Информативные сообщения об ошибках валидации.
\end{enumerate}

\subsection{Специальные операции}

\begin{enumerate}[leftmargin=*]
    \item Количество объектов с \texttt{impactSpeed} равным заданному.
    \item Количество объектов с \texttt{weaponType} меньше заданного.
    \item Количество объектов с \texttt{impactSpeed} больше заданного.
    \item Удалить всех героев без зубочисток.
    \item Пересадить всех героев без машины на красные «Lada Kalina».
\end{enumerate}
Спецоперации реализуются в бизнес-логике без прямого использования функций и процедур БД.

\subsection{Требования к хранению и технологиям}

\begin{itemize}[leftmargin=*]
    \item СУБД: PostgreSQL, БД \texttt{studs}, хост \texttt{pg}.
    \item Все ограничения полей — на уровне ORM и БД. ID генерируется средствами БД.
    \item Меню для навигации между логическими частями системы.
    \item Стек: Spring Boot, JPA + Hibernate, разделение на слои.
\end{itemize}

\newpage
%============================================================
\section{UML-диаграммы}

\subsection{Диаграмма классов}

\begin{figure}[h!]
\centering
\begin{tikzpicture}
  \begin{class}[text width=5.5cm]{HumanBeing}{0,0}
    \attribute{- id : Integer}
    \attribute{- name : String}
    \attribute{- coordinates : Coordinates}
    \attribute{- creationDate : LocalDateTime}
    \attribute{- realHero : boolean}
    \attribute{- hasToothpick : boolean}
    \attribute{- car : Car}
    \attribute{- mood : Mood}
    \attribute{- impactSpeed : float}
    \attribute{- weaponType : WeaponType}
  \end{class}

  \begin{class}[text width=3.5cm]{Coordinates}{-6,-7}
    \attribute{- id : Long}
    \attribute{- x : Float}
    \attribute{- y : float}
  \end{class}

  \begin{class}[text width=2.5cm]{Car}{3,-7}
    \attribute{- id : Long}
    \attribute{- name : String}
  \end{class}

  \begin{class}[text width=2.5cm]{Mood}{-9,-3}
    \operation{<<enum>>}
    \attribute{SADNESS}
    \attribute{APATHY}
    \attribute{CALM}
    \attribute{FRENZY}
  \end{class}

  \begin{class}[text width=2.5cm]{WeaponType}{6,-3}
    \operation{<<enum>>}
    \attribute{PISTOL}
    \attribute{KNIFE}
    \attribute{BAT}
  \end{class}

  \draw[umlcd style, ->] (HumanBeing) -- node[above,sloped] {1} node[below,sloped] {*} (Coordinates);
  \draw[umlcd style, ->] (HumanBeing) -- node[above,sloped] {0..1} node[below,sloped] {*} (Car);
  \draw[umlcd style, ->, dashed] (HumanBeing) -- (Mood);
  \draw[umlcd style, ->, dashed] (HumanBeing) -- (WeaponType);
\end{tikzpicture}
\caption{Диаграмма классов предметной области}
\end{figure}

\newpage
\subsection{Диаграмма пакетов}

\begin{figure}[h!]
\centering
\begin{tikzpicture}[node distance=1.2cm, auto,
  pkg/.style={rectangle, draw, minimum width=4cm, minimum height=0.8cm, font=\ttfamily\small},
  layer/.style={rectangle, draw, dashed, minimum width=13cm, minimum height=2.8cm, fill=gray!5, label={[font=\bfseries]above:#1}}]

  % Layers
  \node[layer=Presentation] (pres) at (0,0) {};

  \node[layer=Business] (bus) at (0,-4) {};

  \node[layer=Persistence] (per) at (0,-8) {};

  % Packages
  \node[pkg] (ctrl) at (-4,0.3) {controller};
  \node[pkg] (ws) at (4,0.3) {websocket};
  \node[pkg] (dto) at (0,0.3) {dto};

  \node[pkg] (svc) at (-3.5,-3.7) {service};
  \node[pkg] (spec) at (1,-3.7) {specification};
  \node[pkg] (map) at (5,-3.7) {mapper};

  \node[pkg] (repo) at (-2,-7.7) {repository};
  \node[pkg] (ent) at (3,-7.7) {entity};

  % Arrows between layers
  \draw[->, thick] (-6,-1.2) -- node[right] {} (-6,-2.8);
  \draw[->, thick] (-6,-5.2) -- node[right] {} (-6,-6.8);

\end{tikzpicture}
\caption{Диаграмма пакетов приложения}
\end{figure}

\vspace{0.5cm}

\textbf{Описание пакетов:}

\begin{itemize}[leftmargin=*]
    \item \texttt{controller} — REST-контроллеры (\texttt{HumanBeingController}, \texttt{ReferenceController}, \texttt{SpecialOperationsController}).
    \item \texttt{dto/request} — входные DTO с Bean Validation (\texttt{HumanBeingRequest}, \texttt{CoordinatesRequest}, \texttt{CarRequest}).
    \item \texttt{service} — бизнес-логика (\texttt{HumanBeingService}, \texttt{CoordinatesService}, \texttt{CarService}).
    \item \texttt{repository} — Spring Data JPA репозитории.
    \item \texttt{entity} — JPA-сущности и перечисления (\texttt{HumanBeing}, \texttt{Coordinates}, \texttt{Car}, \texttt{Mood}, \texttt{WeaponType}).
    \item \texttt{mapper} — преобразование DTO в сущности.
    \item \texttt{specification} — динамические фильтры через \texttt{JpaSpecificationExecutor}.
    \item \texttt{websocket} — STOMP-брокер, события \texttt{EntityChangedEvent}, transactional listeners.
    \item \texttt{config} — конфигурация WebSocket.
    \item \texttt{exception} — глобальный обработчик исключений, кастомные исключения.
\end{itemize}

\newpage
%============================================================
\section{Исходный код}

Исходный код проекта доступен в репозитории GitHub:

\url{https://github.com/yaart/IS_lab1}

\subsection{Стек технологий}
\begin{itemize}
    \item \textbf{Backend:} Java 17, Spring Boot 3.4.4, Spring Data JPA, Hibernate, Spring WebSocket
    \item \textbf{Frontend:} React 18, Vite, Material UI (MUI), STOMP over WebSocket
    \item \textbf{База данных:} PostgreSQL 18 (docker-compose), хост \texttt{pg}, БД \texttt{studs}
    \item \textbf{Сборка:} Maven (backend), Vite (frontend)
\end{itemize}

\subsection{Структура проекта}
\begin{lstlisting}[language=,basicstyle=\footnotesize\ttfamily]
is-lab1-spring/
  pom.xml
  docker-compose.yml
  src/main/java/com/example/islab1/
    IsLab1Application.java
    config/WebSocketConfig.java
    controller/
      SpaController.java
      api/
        HumanBeingController.java
        ReferenceController.java
        SpecialOperationsController.java
    dto/request/
      CarRequest.java
      CoordinatesRequest.java
      HumanBeingRequest.java
    entity/
      Car.java, Coordinates.java, HumanBeing.java
      Mood.java, WeaponType.java
    exception/
      ApiErrorResponse.java, BusinessException.java
      GlobalExceptionHandler.java, ResourceNotFoundException.java
      ValidationErrorResponse.java
    mapper/HumanBeingMapper.java
    repository/
      CarRepository.java, CoordinatesRepository.java
      HumanBeingRepository.java
    service/
      CarService.java, CoordinatesService.java
      HumanBeingService.java
    specification/HumanBeingSpecification.java
    websocket/
      ChangeType.java, EntityChangedEvent.java
      EntityChangedListener.java, EntityChangePublisher.java
  src/main/resources/
    application.yml, application-dev.yml
  frontend/
    src/App.jsx
    src/components/
      HumansTable.jsx, HumanDialog.jsx
      SpecialOps.jsx, ConfirmDialog.jsx
    src/api.js, src/websocket.js, src/theme.js
\end{lstlisting}

\subsection{Схема базы данных}
\begin{lstlisting}[language=SQL]
CREATE TABLE coordinates (
    id BIGSERIAL PRIMARY KEY,
    x FLOAT NOT NULL CHECK (x <= 256),
    y FLOAT NOT NULL CHECK (y > -533)
);

CREATE TABLE car (
    id BIGSERIAL PRIMARY KEY,
    name VARCHAR(255)
);

CREATE TABLE human_being (
    id SERIAL PRIMARY KEY,
    name VARCHAR(255) NOT NULL CHECK (char_length(trim(name)) > 0),
    coordinates_id BIGINT NOT NULL REFERENCES coordinates(id),
    creation_date TIMESTAMP NOT NULL,
    real_hero BOOLEAN NOT NULL DEFAULT FALSE,
    has_toothpick BOOLEAN NOT NULL DEFAULT FALSE,
    car_id BIGINT REFERENCES car(id),
    mood VARCHAR(20),
    impact_speed FLOAT CHECK (impact_speed <= 738),
    weapon_type VARCHAR(20) NOT NULL
);
\end{lstlisting}

\newpage
%============================================================
\section{Выводы}

В ходе выполнения лабораторной работы была разработана информационная система для управления объектами класса \texttt{HumanBeing}. Система реализует полный набор CRUD-операций, фильтрацию и сортировку таблицы, real-time синхронизацию между клиентами через WebSocket, пять специальных операций и защиту от некорректного удаления связанных объектов.

Применён трёхслойный архитектурный подход (контроллеры, сервисы, репозитории) с чётким разделением ответственности. Все ограничения полей классов дублированы на уровнях Java-валидации, ORM и базы данных, что обеспечивает целостность данных на всех уровнях приложения.

Для клиентской части использован стек React + Material UI с серверной пагинацией и мгновенным обновлением данных через STOMP-протокол. Навигация между таблицей объектов и специальными операциями реализована через меню с вкладками.

\end{document}